A Library represents a single shared library (frequently
referred to as a DLL or DSO, depending on the OS) that has been loaded into the
target process. In addition, a Library will be used to represent the
process\' executable. Process\' with statically linked executables
will only contain the single Library that
represents the executable.

Each Library contains a load address and a file
name. The load address is the address at which the OS loaded the library, and
the file name is the path to the library\'s file. Note
that on some operating systems (e.g., Linux and FreeBSD) the load
address does not necessarily represent the beginning of the library in memory;
instead it is a value that can be added to a library\'s
symbol offsets to compute the dynamic address of a symbol.

Libraries may be loaded and unloaded by the process during execution. A
library load or unload can trigger a callback with an EventLibrary parameter.
The current list of libraries loaded into a process can be accessed via the
\code{LibraryPool} of a \code{Process}.

\begin{apient}
  Library::ptr
  Library::const_ptr
\end{apient}

\apidesc{
  \code{typedef}s for a pointer and a const pointer to a library.
}

\begin{apient}
  std::string getName() const
\end{apient}

\apidesc{
  Returns the file name for this Library.
}

\begin{apient}
  std::string getAbsoluteName() const
\end{apient}

\apidesc{
  Returns a file name for this Library that does not contain symlinks or a
  relative path.
}

\begin{apient}
  Dyninst::Address getLoadAddress() const
\end{apient}

\apidesc{
  Returns the load address for this Library.
}

\begin{apient}
  Dyninst::Address getDataLoadAddress() const
\end{apient}

\apidesc{
  The AIX operating system can have two load addresses for a library:
  one for the
  code region and one for the data region. On AIX Library::getLoadAddress
  returns the load address of the code region and Library::getDataLoadAddress
  returns the load address of the data region. On non-AIX systems this function
  returns 0.
}

\begin{apient}
  Dyninst::Address getDynamicAddress() const
\end{apient}

\apidesc{
  On ELF based systems (FreeBSD, Linux) this function will return the
  address of the Library\'s dynamic section. On other
  systems this function returns 0.
}

\begin{apient}
  bool isSharedLib() const
\end{apient}

\apidesc{
  Returns true of this Library object is refers to a shared library
  and false if it refers to an executable.
}

\begin{apient}
  void *getData() const
\end{apient}

\apidesc{
  Returns an opaque data object that user code can associate with
  this library. Use setData to set this opaque value.
}

\begin{apient}
  void setData(void *p) const
\end{apient}

\apidesc{
  This function sets associates an opaque data object with the Library.
  ProcControlAPI does not try to interpret this value, but will return it with
  the getData function.
}
